
%(BEGIN_QUESTION)
% Copyright 2006, Tony R. Kuphaldt, released under the Creative Commons Attribution License (v 1.0)
% This means you may do almost anything with this work of mine, so long as you give me proper credit

Ledningsevneceller oppgis med en {\it cellekonstant}, ofte skrevet med gresk bokstav $\theta$ i ledningsevneligninger. Riktig enhet for cellekonstanten er inverse centimeter, cm$^{-1}$.

Ta følgende to ligninger og kombiner dem ved algebraisk substitusjon for å få en ny ligning som løser væskens spesifikke ledningsevne ($k$) ut fra rå ledningsevne ($G$) og cellekonstant ($\theta$):

$$G = k{A \over d} \hbox{\hskip 100pt} \theta = {d \over A}$$

\noindent
Der,

$G$ = Ledningsevne i Siemens (S)

$\theta$ = Cellekonstant i cm$^{-1}$

$k$ = Spesifikk ledningsevne i Siemens per centimeter (S/cm)

$A$ = Elektrodeareal i cm$^{2}$

$d$ = Elektrodeavstand i cm

\vskip 10pt

\underbar{file i00607}
%(END_QUESTION)





%(BEGIN_ANSWER)

$$k = G \theta$$

%(END_ANSWER)





%(BEGIN_NOTES)


%INDEX% Måling, analytisk: ledningsevne
%INDEX% Matematikkrepetisjon: manipulering og kombinasjon av ligninger til ny ligning

%(END_NOTES)

